\chapter{Reactive Arthritis}
\index{Reactive Arthritis}
\label{sec:Reactive Arthritis}

\section{Overview}
Reactive arthritis is an aseptic inflammatory arthritis occurring 1--4 weeks after a genitourinary (usually \textit{Chlamydia trachomatis}) or digestive tract (usually \textit{Salmonella}, \textit{Shigella}, \textit{Yersinia}, \textit{Campylobacter}) infection. This autoimmune condition occurs when the immune system mistakenly attacks the body's own tissues. It follows a chronic course with periods of flare and remission.
\section{Causes}
The underlying mechanisms involve complex interactions between genetic predisposition and environmental factors. Disruption of normal physiological processes leads to the characteristic features of the condition. Multiple contributing factors may act synergistically to trigger or worsen the condition.

\begin{itemize}[noitemsep]
  \item This condition happens because of an immune response to the initial infection triggering joint inflammation. The classic triad is arthritis, urethritis, and conjunctivitis ("can't see, can't pee, can't climb a tree"). Arthritis is typically asymmetric and oligoarticular, affecting large joints of the lower extremities
  \item Enthesitis, dactylitis, and oral ulcers are common. Most cases resolve within 3--12 months.
\end{itemize}
\section{Risk Factors}


\begin{itemize}[noitemsep]
  \item Genetic predisposition and family history
  \item Advancing age
  \item Lifestyle factors (diet, exercise, smoking, alcohol)
  \item Environmental exposures
  \item Underlying medical conditions
  \item Medication use
\end{itemize}
\section{Signs and Symptoms}

\subsection{Common symptoms}
\begin{itemize}[noitemsep]
  \item Symptoms vary depending on the severity and extent of the condition Pain or discomfort in the affected area Functional impairment affecting daily activities Changes in appearance or function of affected tissues Systemic symptoms may include fatigue, fever, or weight changes Symptoms may develop gradually or appear suddenly
  \item Pain or discomfort in the affected area
  \end{itemize}

\subsection{Emergency warning signs}
\begin{itemize}[noitemsep]
  \item Severe or worsening symptoms of reactive arthritis require immediate medical evaluation
\end{itemize}
\section{Progression and Stages}
It is strongly associated with HLA-B27 and classified as a seronegative spondyloarthropathy. The condition may progress through different stages of severity over time. Early stages often involve mild or subtle symptoms that may be overlooked. Moderate stages involve more noticeable symptoms affecting daily function. Advanced stages may involve significant impairment and complications.
\section{Complications}
\begin{itemize}[noitemsep]
  \item Disease progression leading to worsening symptoms
  \item Secondary infections or injuries
  \item Side effects of treatment
  \item Reduced quality of life
  \item Emotional and psychological impact
\end{itemize}
\section{Diagnosis}
Doctors usually diagnose this based on your symptoms and a physical exam with history of preceding infection. Diagnosis is based on a combination of clinical features, laboratory findings (autoantibodies, inflammatory markers), and imaging studies. Classification criteria help standardize the diagnostic process. Comprehensive medical history and physical examination Laboratory tests to assess organ function and rule out other conditions Imaging studies to visualize affected structures Specialized testing depending on the affected system Consultation with relevant specialists Monitoring of disease progression over time

\begin{itemize}[noitemsep]
  \item Comprehensive medical history and physical examination
  \item Laboratory tests to assess organ function
  \item Imaging studies to visualize affected structures
  \item Specialized testing depending on the affected system
  \item Consultation with relevant specialists
\end{itemize}
\section{Prevention}
\begin{itemize}[noitemsep]
  \item Maintain a healthy lifestyle with balanced diet and exercise
  \item Avoid known risk factors
  \item Attend regular medical check-ups for early detection
  \item Follow recommended screening guidelines
\end{itemize}
\section{Treatment Options}
\begin{itemize}[noitemsep]
  \item Treatment options include NSAIDs, intra-articular corticosteroids, and for chronic or severe disease, DMARDs (sulfasalazine) or TNF inhibitors
  \item Antibiotics are not helpful once arthritis has developed.
  \item Medications to manage symptoms and address underlying mechanisms
  \item Lifestyle modifications and self-care measures
  \item Physical therapy and rehabilitation
  \item Surgical intervention when indicated
  \item Regular monitoring and follow-up care
\end{itemize}
\section{Living with It}
\begin{itemize}[noitemsep]
  \item Follow your treatment plan as prescribed
  \item Maintain regular follow-up with your healthcare provider
  \item Make necessary lifestyle adjustments
  \item Seek support from family, friends, or support groups
  \item Stay informed about your condition
\end{itemize}
\section{Outlook}
\begin{itemize}[noitemsep]
  \item In most cases, the outlook is good
  \item Most cases resolve within 3--12 months.
\end{itemize}
